% ==================================================
% Forward Difference
% Discrete Calculus — Difference Operators
% ==================================================

\section{Difference Operators}
\label{sec:difference-operators}

Discrete calculus studies functions defined on the integers and the way
their values change from one integer to the next. The central tool for
measuring this change is the \emph{difference operator}, which plays a role
analogous to the derivative in continuous calculus.

\section*{Motivation}

In ordinary calculus the derivative measures the rate of change of a function:
\[
f'(x) = \lim_{h\to 0} \frac{f(x+h)-f(x)}{h}.
\]
In the discrete setting the independent variable changes in steps of size
\(1\). Instead of taking a limit, we measure the difference between
successive values of the function directly.

\section*{Forward Difference}

\begin{tcolorbox}[colback=defbox, colframe=defborder, arc=2pt,
  left=6pt, right=6pt, top=4pt, bottom=4pt,
  title={\small\textbf{Definition (Forward Difference Operator)}},
  fonttitle=\small\bfseries]
Let \(f : \mathbb{Z} \to \mathbb{R}\).
The \emph{forward difference operator} \(\Delta\) is defined by
\[
\Delta f(n) = f(n+1) - f(n).
\]
The function \(\Delta f\) measures the change in \(f\) when the input
increases by one unit.
\end{tcolorbox}
\label{def:forward-difference}

\begin{remark*}[Fully quantified form]
\(\forall f : \mathbb{Z}\to\mathbb{R},\; \forall n \in \mathbb{Z}:\;
\Delta f(n) := f(n+1) - f(n)\).
\end{remark*}

\begin{remark*}[Analogy with differentiation]
The operator \(\Delta\) is the discrete analogue of the derivative.
While the derivative measures infinitesimal change through a limit,
the forward difference measures change between successive integer inputs
exactly, with no limit required.
\end{remark*}

\section*{Example}

Consider the function \(f(n) = n^2\). Applying the difference operator gives
\[
\begin{aligned}
\Delta f(n)
&= f(n+1) - f(n) \\
&= (n+1)^2 - n^2 \\
&= 2n + 1.
\end{aligned}
\]
Thus the forward difference of the quadratic function \(n^2\) is the linear
function \(2n+1\).

\begin{remark*}[Degree reduction]
This mirrors a familiar fact from ordinary calculus: \(\tfrac{d}{dx} x^2 = 2x\).
Just as derivatives reduce the degree of a polynomial by one,
finite differences also lower the degree of polynomial sequences.
Note, however, that \(\Delta(n^2) = 2n+1\) rather than \(2n\); ordinary
powers do not behave as cleanly under \(\Delta\) as they do under
differentiation. This motivates introducing the \emph{falling factorial}
basis later in this chapter.
\end{remark*}
